\scanpage{22}
\mspage{13}
% archive sheet number: 22
% lines 1--5
Point de vue ``symétrique''. Structure sur deux ensembles de base $M$ (multistrates), $\unsure{F|\mathscr{F}}$ (figures)

Structure\par
a) Relation d'incidence $X\triangleleft F$ entre $M,F$ (suffira à définir $\leq$ sur $M$ et sur $F$)

b) Relations $\ll$, $\mid\circ\mid$ sur $M$

% lines 6--7
Pour tout $F\in\unsure{\mathscr{F}|\mathcal{F}}$, on pose
\[
\widetilde{F}=\{X\in M\mid X\triangleleft F\}\qquad
\text{On pose }F\leq G\overset{\text{déf}}{\Longleftrightarrow}\widetilde{F}\subset\widetilde{G}
\]

% lines 8--14: entire draft block crossed by long diagonals
\margin{Ainsi,}
\deleted{\uline{At M1} $\widetilde{F}=\widetilde{G}\unsure{\Rightarrow}\Rightarrow F=G$, i.e. $\mathscr{F}\longrightarrow\unsure{\mathscr{P}|\mathcal{P}}(M)$, $F\longmapsto\widetilde{F}$ est \unsure{injection|injective}, ou encore le préordre $\leq$ sur $\mathscr{F}$ est un \uuline{ordre}.}

\deleted{\uline{At M2} $\unsure{F}\ \widetilde{\mathscr{F}}$ \unsure{couvre} est une \unsure{figure élémentaire dans}\inserted{famille de parties de $M$ quasi-admissible} $\unsure{\mathscr{F}}$ qui couvre $M$ et sépare les points de $M$, i.e.}

\deleted{a) $\forall X\in M$, $\unsure{\mathscr{F}}\ \mathscr{F}^{X}=\{F\in\mathscr{F}\mid X\triangleleft F\}$ a un plus petit élément (on l'appelle la \uline{figure élémentaire} ou multistrate déterminée par $X$)}

\deleted{\unsure{b)}}

% lines 15--20: revised axioms
\uline{At M1} La relation $X\triangleleft F$ est séparante en $M$ et en $\mathscr{F}$, i.e.

a) L'application $F\longmapsto\widetilde{F}$, $\mathscr{F}\longrightarrow\mathscr{P}(M)$ est injective [ou encore, $\leq$ est une relation d'\uline{ordre}]

b) $\deleted{\unsure{\forall}}$ Si $X,Y\in M$ tels que $X\neq Y$, alors $\exists F\in\mathscr{F}$ avec $X\in\widetilde{F}$, $Y\notin\widetilde{F}$ ou l'inverse (i.e. $\widetilde{\mathscr{F}}\subset\mathscr{P}(M)$ sépare les points de $M$)

% lines 21--24
\uline{At M2} La famille de parties $\widetilde{\mathscr{F}}$ dans $M$ est quasi-admissible, i.e. $\forall X\in M$, $\deleted{\unsure{\mathscr{F}}}$ posant $\mathscr{F}^{X}=\{F\in\mathscr{F}\mid X\triangleleft F\}$, $\mathscr{F}^{X}$ a un plus petit élément (appelé la \uline{figure élémentaire} associée à la multistrate $X$)

% lines 25--29
Ainsi \deleted{$X\triangleleft F$ équivaut à}
\[
X\triangleleft F\Longleftrightarrow F_X\leq F\qquad(\Longleftrightarrow X\in\widetilde{F})
\]
de sorte que la condition At M1 b), compte tenu de At M1 a), équivaut à
\[
M\longrightarrow\mathscr{F}\quad X\longrightarrow F_X\qquad\text{est injectif}.
\]

% lines 30--32
\uline{At M3} a) Soit $\deleted{\unsure{\mathscr{F}}}\ \unsure{\underline{\Phi}_i|\underline{\varphi}_i},\underline{\Phi}\in\widetilde{\mathscr{F}}$, avec $\underline{\Phi}_i\subset\underline{\Phi}$. Alors $\bigcup\underline{\Phi}_i\in\widetilde{\mathscr{F}}$ (stabilité de $\widetilde{\mathscr{F}}$ par $\unsure{\bigcup_{\alpha}|\bigcup}$ \unsure{majorées})

\margin{résulte du reste si $M\neq\varnothing$}
b) $\unsure{\widetilde{\mathscr{F}}|\mathscr{F}}\neq\varnothing$ ($\deleted{\Longleftrightarrow}\varnothing\in\widetilde{\mathscr{F}}\Longleftrightarrow\widetilde{\mathscr{F}}\neq\varnothing\Longleftrightarrow\mathscr{F}$ a un plus petit élément)
